%%================================================
%% Chapter 5 - สรุปผล ข้อจำกัด และข้อเสนอแนะ
%%================================================
\chapter{สรุปผล ข้อจำกัด และข้อเสนอแนะ}
\label{chapter5}

\section{สรุปภาพรวมของระบบ}
\label{sec:ch5_system_summary}

งานวิจัยนี้พัฒนาระบบต้นแบบ WheelSense สำหรับสภาพแวดล้อมอัจฉริยะที่ช่วยสนับสนุนการดูแลผู้ใช้เก้าอี้รถเข็นและผู้สูงอายุ โดยบูรณาการข้อมูลจาก sensor บนเก้าอี้รถเข็น อุปกรณ์สวมใส่, mobile gateway, server, dashboard และระบบ AI/MCP ให้อยู่ในโครงสร้างเดียวกัน ระบบที่พัฒนาขึ้นใช้ M5StickC Plus2 เป็น BLE wheelchair sensor สำหรับอ่านข้อมูล IMU และ battery ใช้ Polar Verity Sense สำหรับข้อมูลสัญญาณชีพ ใช้แอปพลิเคชัน Flutter บนโทรศัพท์เป็น gateway หลัก และส่งข้อมูลเข้าสู่ backend ผ่าน MQTT หรือ REST API จากนั้นข้อมูลถูกจัดเก็บและให้บริการผ่าน FastAPI, PostgreSQL/JSONB, MQTT broker, Next.js dashboard และ AI/MCP pipeline ที่ออกแบบให้มีการจำกัด schema และตรวจสอบก่อนสั่งงาน

ผลการประเมินในบทที่~\ref{chapter4} สนับสนุนว่าระบบสามารถทำงานร่วมกันได้ในระดับต้นแบบ โดยเฉพาะส่วนการรับส่งข้อมูลจากอุปกรณ์ภาคสนามและ mobile gateway การทดสอบ M5StickC Plus2 พบว่าอุปกรณ์อ่านข้อมูล IMU ได้ที่ค่าเฉลี่ย 19.6 Hz ใกล้เคียงเป้าหมาย 20 Hz และมีความครบถ้วนของ telemetry เท่ากับ 98.4\% ส่วนการจำแนกสถานะการเคลื่อนไหวจากข้อมูล IMU ให้ค่าความถูกต้องในช่วง 0.90--0.98 ตามกลุ่มสถานะที่ทดสอบ ผลดังกล่าวแสดงว่าการใช้ gyro offset, EMA, deadband และการคำนวณค่าจาก IMU สามารถใช้เป็นพื้นฐานสำหรับการติดตามสถานะการเคลื่อนไหวของรถเข็นในงานต้นแบบได้

ในส่วน mobile gateway และ server ระบบสามารถรวมข้อมูลจาก Polar Verity Sense, M5StickC Plus2 และข้อมูลจากโทรศัพท์ก่อนส่งขึ้น platform ได้ การทดสอบพบว่า BLE GATT connection มีเวลาเชื่อมต่อ p50 เท่ากับ 3.2 วินาที และ p95 เท่ากับ 5.8 วินาที ขณะที่เวลาแฝงจาก smartphone ไปถึง broker มีค่า p50 เท่ากับ 78 ms และ p95 เท่ากับ 215 ms อัตราความสำเร็จของ MQTT publish เท่ากับ 98.6\% และความครบถ้วนของ vitals telemetry อยู่ในช่วง 96.2--98.9\% ส่วน server และ database มีเวลาแฝงในระดับที่เพียงพอสำหรับ dashboard ต้นแบบ เช่น การอ่านสถานะล่าสุดของผู้พักอาศัยมีค่า p95 เท่ากับ 43 ms และการอ่าน audit trail มีค่า p95 เท่ากับ 51 ms

สำหรับการระบุตำแหน่งภายในอาคาร การทดสอบในบทที่~\ref{chapter4} ใช้ข้อมูล RSSI fingerprint และ KNN เพื่อประเมินแนวทางการจำแนกโซน ผลที่มีหลักฐานรองรับในปัจจุบันเป็นผลรายกรณีของชุดทดสอบ โดยมีค่าความถูกต้องและ macro-F1 ในกรณีที่ 1 เท่ากับ 0.8857 และ 0.8889 กรณีที่ 2 เท่ากับ 0.6857 และ 0.6667 และกรณีที่ 3 เท่ากับ 0.8857 และ 0.8857 ตามลำดับ ผลนี้สะท้อนว่า KNN สามารถใช้เป็นแนวทางพื้นฐานสำหรับการบอกตำแหน่งระดับโซนได้ แต่ค่าความถูกต้องยังขึ้นอยู่กับสภาพแวดล้อมและความเสถียรของ RSSI จึงไม่ควรตีความเป็นค่าความแม่นยำรวมของอาคารทั้งหมดจนกว่าจะมีชุดข้อมูลดิบที่ครอบคลุมกว่าเดิม

ในส่วน dashboard และ HMI ผลการประเมินชี้ว่าผู้ใช้สามารถทำงานหลักบนหน้าจอระบบได้ในระดับต้นแบบ โดย task completion เฉลี่ยเท่ากับ 0.93 คะแนนความพึงพอใจเฉลี่ยเท่ากับ 4.31/5 คะแนนความอ่านง่ายของ alert เท่ากับ 4.42/5 และคะแนนการนำทางตามบทบาทผู้ใช้เท่ากับ 4.18/5 ผลดังกล่าวสะท้อนว่า dashboard ช่วยให้ผู้ดูแลเห็นข้อมูลผู้พักอาศัย สถานะอุปกรณ์ การแจ้งเตือน และบริบทการดูแลจากจุดศูนย์กลางได้ อย่างไรก็ตามค่า SUS mean 76.5 และ SD 8.7 ที่พบในเอกสารประกอบยังต้องแนบ raw item score ของแบบสอบถามให้ครบก่อน จึงควรใช้เป็นผลเบื้องต้นด้าน usability และไม่สรุปเป็นผลยืนยันสุดท้ายของงานวิจัยในตอนนี้

ระบบ AI/MCP pipeline ที่พัฒนาขึ้นมีเป้าหมายเพื่อลดปัญหาของการใช้ LLM แบบรวมศูนย์เดิมที่มี latency สูง ใช้ token มาก และเสี่ยงต่อการสั่งงานผิดบริบท โครงสร้างใหม่แบ่งงานออกเป็น 5 ชั้น ได้แก่ Intent Router, Context Assembler, Behavioral State, Constrained LLM และ Execution Gate โดย Layer 3 ทำงานเป็นข้อมูลบริบทเชิงพฤติกรรม ไม่ได้อยู่ใน execute path หลัก การใช้ schema-bound prompting ร่วมกับ Execution Gate ช่วยให้คำสั่งที่เกี่ยวข้องกับ tool ถูกตรวจสอบตาม schema และ policy ก่อนดำเนินการ ผลการทดสอบ Inter-Rater Reliability มีค่า 31/36 หรือ 86.11\% และข้อผิดพลาดทั้ง 5 กรณีอยู่ในกลุ่มคำสั่งที่มีความกำกวมสูง ซึ่งสอดคล้องกับแนวทางของระบบที่เลือกให้ถามกลับเมื่อความมั่นใจไม่เพียงพอ แทนที่จะเดาคำสั่งแล้วส่งต่อไปยัง tool

ผลการทดสอบ Text-Distance หรือ cosine similarity แสดงว่าคำถามทั่วไปมีค่าเฉลี่ย 0.85 และคำถามด้านสุขภาพมีค่าเฉลี่ย 0.76 ซึ่งสะท้อนว่า RAG และ retrieval logic ช่วยดึงบริบทที่เกี่ยวข้องให้ LLM ใช้ตอบคำถามได้ แต่ลำดับของบริบทและการจัด context window ยังทำให้คะแนนมีความผันผวนได้ สำหรับ latency analysis พบว่าเวลารวมของ pipeline มีค่า p50 เท่ากับ 14,592.43 ms หรือประมาณ 14.59 วินาที และ p95 เท่ากับ 36,941.06 ms หรือประมาณ 36.94 วินาที เมื่อนำไปเทียบกับบันทึกการพัฒนาแนวทางเดิมที่บางรอบใช้เวลาประมาณ 30--45 นาทีต่อรอบการทำงาน จึงเห็นได้ว่าการแยก logic ออกจาก LLM ช่วยลดเวลารอได้มาก แต่ Layer 4 Constrained LLM ยังเป็น bottleneck หลักของระบบ

โดยสรุป ระบบต้นแบบ WheelSense แสดงให้เห็นความเป็นไปได้ของการรวมข้อมูลจาก wheelchair sensor, wearable sensor, mobile gateway, server, dashboard และ AI/MCP pipeline เพื่อช่วยผู้ดูแลติดตามสถานะของผู้พักอาศัยในรูปแบบรวมศูนย์ ระบบไม่ได้มีวัตถุประสงค์เพื่อใช้แทนการประเมินทางการแพทย์หรือแทนอุปกรณ์ทางการแพทย์ แต่เป็นต้นแบบสำหรับการเฝ้าระวังเบื้องต้น การแสดงข้อมูลบริบท และการเสนอ action ให้ผู้ดูแลตรวจสอบก่อนดำเนินการ

\section{ข้อจำกัดของระบบ}
\label{sec:ch5_limitations}

แม้ผลการทดสอบจะแสดงให้เห็นว่าระบบสามารถทำงานได้ในระดับต้นแบบ แต่งานวิจัยนี้ยังมีข้อจำกัดหลายด้านที่ต้องพิจารณาก่อนพัฒนาต่อในสภาพแวดล้อมที่ซับซ้อนกว่าเดิม ข้อจำกัดเหล่านี้ไม่ได้เป็นความล้มเหลวของระบบ แต่เป็นขอบเขตที่งานวิจัยฉบับนี้ยังไม่ได้แก้ครบถ้วน

\begin{enumerate}
    \item \textbf{ข้อจำกัดของ AI routing} ระบบยังอยู่บนแนวคิด single-intent routing เป็นหลัก หากผู้ใช้ส่งคำสั่งที่มีหลายเจตนาพร้อมกัน เช่น ขอข้อมูลสุขภาพ ขอเปลี่ยนสถานะอุปกรณ์ และขอค้นประวัติในคำสั่งเดียวกัน ระบบยังต้องแยกเจตนาให้ชัดเจนก่อนดำเนินการ จึงอาจต้องถามกลับหรือใช้เวลามากขึ้น

    \item \textbf{ข้อจำกัดจากการใช้โมเดลภาษาเดียวกับทุกงาน} การใช้ Gemma 4B เป็นโมเดลหลักสำหรับงานหลายประเภททำให้คำสั่งง่ายบางประเภทใช้เวลามากเกินจำเป็น โดยเฉพาะเมื่อคำสั่งนั้นควรตอบได้ด้วย rule, deterministic matching หรือ model ขนาดเล็กกว่า ผล latency ในบทที่~\ref{chapter4} แสดงชัดว่า Layer 4 เป็นคอขวดหลัก โดยมีค่า p95 เท่ากับ 36.49 วินาทีเฉพาะส่วน LLM

    \item \textbf{ยังไม่มี feedback loop จาก log การใช้งาน} ระบบสามารถเก็บบริบทและผลการสั่งงานได้ แต่ยังไม่มี mechanism ที่นำ log ของคำถาม คำตอบ การแก้ไขของผู้ดูแล หรือกรณีที่ระบบถามกลับ ไปปรับปรุง retrieval, intent routing หรือ prompt template โดยอัตโนมัติ จึงยังต้องปรับปรุงด้วยการวิเคราะห์จากผู้พัฒนาเป็นหลัก

    \item \textbf{Layer 3 และ Edge AI ยังเป็น proof of concept} ชั้น Behavioral State ช่วยเตรียมบริบทด้านการเคลื่อนไหว สัญญาณชีพ และเวลาให้ระบบ AI ใช้ประกอบการตอบคำถาม แต่ยังไม่ได้ผ่านการประเมินระยะยาวหรือการตรวจสอบเชิงคลินิก จึงควรถือเป็นบริบทสนับสนุน ไม่ใช่โมเดลตัดสินผลทางการแพทย์หรือโมเดลตัดสินความเสี่ยงทางการแพทย์

    \item \textbf{ข้อมูลทดสอบจริงระยะยาวยังไม่เพียงพอ} การทดสอบในงานวิจัยนี้เป็นการทดสอบภายในพื้นที่ทดลองและชุดข้อมูลที่จัดเตรียมขึ้นเพื่อประเมิน subsystem หลัก ยังไม่มีข้อมูลระยะยาวจากผู้ใช้หลายกลุ่ม หลายช่วงเวลา และหลายสภาพแวดล้อม จึงยังไม่เพียงพอสำหรับสรุปความทนทานของระบบในทุกสถานการณ์

    \item \textbf{ระบบยังไม่ใช่อุปกรณ์ทางการแพทย์ที่ผ่านการรับรอง} ข้อมูล HR, PPG, IMU และตำแหน่งที่ระบบแสดงผลเป็นข้อมูลเพื่อสนับสนุนการเฝ้าระวังเบื้องต้นเท่านั้น ไม่ควรถูกใช้แทนการประเมินจากบุคลากรทางการแพทย์ และยังไม่ได้ผ่านกระบวนการรับรองมาตรฐานของอุปกรณ์ทางการแพทย์

    \item \textbf{BLE/RSSI ยังผันผวนตามสภาพแวดล้อม} ผลการทดสอบ localization แสดงว่าความถูกต้องของ KNN ขึ้นอยู่กับรูปแบบ RSSI fingerprint ในแต่ละพื้นที่ เมื่อมีสิ่งกีดขวาง การสะท้อนสัญญาณ หรือการเปลี่ยนตำแหน่งของอุปกรณ์ ค่า RSSI สามารถเปลี่ยนได้และทำให้เกิดความผิดพลาดระหว่างโซนที่อยู่ใกล้กัน

    \item \textbf{หลักฐานด้าน usability และ SUS ยังจำกัด} การประเมิน dashboard/HMI มีจำนวนผู้เข้าร่วมและสถานการณ์ทดสอบในระดับต้นแบบ ข้อมูล task completion และคะแนนความพึงพอใจช่วยให้เห็นแนวโน้มการใช้งาน แต่ค่า SUS 76.5 ยังต้องมี raw item score ครบถ้วนก่อนจึงจะใช้เป็นผลยืนยันได้ นอกจากนี้ยังควรทดสอบซ้ำกับผู้ใช้จริงในบทบาทที่หลากหลายกว่าเดิม
\end{enumerate}

\section{ข้อเสนอแนะในการพัฒนาต่อ}
\label{sec:ch5_future_work}

จากผลการพัฒนาและข้อจำกัดของระบบ แนวทางพัฒนาต่อควรเน้นทั้งด้าน architecture, dataset, AI pipeline, dashboard workflow และการประเมินในสภาพแวดล้อมที่ใกล้เคียงงานดูแลจริงมากขึ้น โดยข้อเสนอแนะหลักมีดังนี้

\begin{enumerate}
    \item \textbf{พัฒนา multi-intent routing} ระบบควรรองรับคำสั่งที่มีหลายเจตนาในประโยคเดียว โดยแยก intent เป็นงานย่อย ตรวจสอบ policy ของแต่ละงาน และจัดลำดับการดำเนินการก่อนส่งไปยัง Layer ถัดไป แนวทางนี้จะช่วยลดการถามกลับในกรณีที่คำสั่งซับซ้อนแต่ยังตีความได้ชัดเจน

    \item \textbf{ใช้ dynamic model multiplexing} ควรแยกประเภทงานตามความซับซ้อน เช่น งานค้นข้อมูลทั่วไป งานอ่านข้อมูล sensor งานสรุปผลสุขภาพเบื้องต้น และงานที่ต้องใช้ tool execution แล้วเลือกใช้ rule, model ขนาดเล็ก หรือ LLM ขนาดใหญ่ตามความเหมาะสม วิธีนี้จะช่วยลด latency ของคำสั่งง่ายและลดภาระของ Layer 4

    \item \textbf{เพิ่ม RAG-based feedback loop} ระบบควรเก็บตัวอย่างคำถาม คำตอบที่ผู้ดูแลยืนยัน คำตอบที่ถูกแก้ไข และกรณีที่ระบบถามกลับ เพื่อนำไปสร้าง feedback dataset สำหรับปรับ ranking, context packing และตัวอย่าง prompt ใน retrieval logic การทำเช่นนี้จะช่วยให้ระบบตอบคำถามเฉพาะบริบทของสถานดูแลได้แม่นยำขึ้น

    \item \textbf{เก็บ dataset ระยะยาวจากหลายสภาพแวดล้อม} ควรเก็บข้อมูล IMU, RSSI, HR, PPG, step counter และ event log ในช่วงเวลาที่ยาวขึ้น รวมถึงทดสอบหลายพื้นที่ หลายรูปแบบการเคลื่อนที่ และหลายตำแหน่ง beacon เพื่อประเมินความเสถียรของระบบและลด bias จากพื้นที่ทดลองเดียว

    \item \textbf{ปรับ Edge AI และ Behavioral State ให้เหมาะกับระบบระยะยาว} Layer 3 ควรถูกพัฒนาด้วยเกณฑ์ที่ตรวจสอบได้ เช่น threshold ของความผิดปกติ ความต่อเนื่องของสถานะ และการจัดลำดับความเร่งด่วนของเหตุการณ์ พร้อมทั้งต้องแยกชัดเจนระหว่างการเฝ้าระวังเบื้องต้นกับการประเมินทางการแพทย์

    \item \textbf{เพิ่ม workflow และ alert management บน dashboard} Dashboard ควรมี queue ของ alert การยืนยันการรับทราบ การบันทึกผู้รับผิดชอบ ประวัติการแก้ไขเหตุการณ์ การกรองตาม role และรายงานสรุปเป็นช่วงเวลา เพื่อให้ผู้ดูแลใช้ระบบได้ต่อเนื่องในงานประจำวันมากขึ้น

    \item \textbf{ปรับปรุง mobile gateway ให้ทนต่อสภาพเครือข่ายจริง} แอป gateway ควรมี offline buffer, reconnect strategy, device diagnostic, battery/status indicator และกลไกแจ้งเตือนเมื่อ Polar หรือ M5StickC Plus2 หลุดการเชื่อมต่อ เพื่อให้การรับส่งข้อมูลไม่ขาดตอนง่ายเมื่ออยู่ในอาคารที่สัญญาณไม่คงที่

    \item \textbf{ขยายการประเมิน usability และ SUS} ควรเก็บ raw SUS item score ให้ครบ เพิ่มจำนวนผู้เข้าร่วม แยกกลุ่มตามบทบาท เช่น Admin, Head Nurse, Supervisor, Observer และ Patient รวมถึงให้ผู้ใช้ทำภารกิจซ้ำหลายรอบ เพื่อวัด learning effect และความสม่ำเสมอของการใช้งาน dashboard

    \item \textbf{ทดสอบในสภาพแวดล้อมที่สมจริงมากขึ้น} การทดสอบรอบต่อไปควรใช้พื้นที่ที่มีรูปแบบการเคลื่อนที่และเครือข่ายใกล้เคียงสถานดูแลผู้สูงอายุจริงมากขึ้น พร้อมกำหนดขั้นตอนด้านความเป็นส่วนตัว การขอความยินยอม และขอบเขตการใช้ข้อมูลอย่างชัดเจน เพื่อให้ผลการประเมินสะท้อนข้อจำกัดของระบบได้ครบกว่าเดิม
\end{enumerate}

ข้อเสนอแนะเหล่านี้สามารถใช้เป็นแนวทางต่อยอดระบบ WheelSense จากต้นแบบที่ทดสอบ subsystem หลักแล้ว ไปสู่ระบบที่มีความเสถียรด้านข้อมูลมากขึ้น มี AI ที่ตอบสนองเร็วขึ้น และมี workflow ที่เหมาะกับการใช้งานของผู้ดูแลในระยะยาว ทั้งนี้การพัฒนาต่อควรคงหลักการสำคัญของงานวิจัยนี้ไว้ คือให้ AI เป็นเครื่องมือช่วยสรุปและเสนอแนวทาง โดยยังให้ผู้ดูแลเป็นผู้ตรวจสอบและตัดสินใจก่อนดำเนินการเสมอ
